\section{The primitive half-class and the limits of a norm invariant}
\label{sec:primitive-weil}

Let $t_0=t_*/2$ on the fixed polarized product $A=E^4=B\times B$.
The lower bound for $t_*$ does not directly give the lower bound for
$t_0$: its Weil coordinate has modulus $\sqrt3$, so a degree bound alone
leaves a possible representation with multiplier one. The odd K\"unneth
components exclude that possibility.

\begin{lemma}[V01: integrality and primitiveness]
The class $t_0$ is integral and primitive in $H^4(A,\Z(2))$, and
$t_0^2=6$.
\end{lemma}

\begin{proof}
The integral K\"unneth decomposition has no torsion. On its summand with
second-factor degree $j$, pullback by $\id\times[-1]$ acts as $(-1)^j$.
Consequently
$([\Gamma_a]+[\Gamma_{-a}])/2$ is the sum of the even K\"unneth
components of the integral class $[\Gamma_a]$, and is integral.
Their difference for $a=1,\omega$ is $t_0$.

Each graph has self-intersection zero. For $a\ne b$ in the four-element
list $1,-1,\omega,-\omega$, its intersection with $\Gamma_b$ is
$\operatorname{Nm}(a-b)^2$. The six contributions to half of $t_*^2$
are $16,16,-9,-9,-1,-1$. Thus $t_*^2=24$ and $t_0^2=6$.
If $t_0=n v$ for an integral class and $n>1$, then $n^2$ would divide $6$,
which is impossible. Hence $t_0$ is primitive.
\end{proof}

\begin{lemma}[V02: small surfaces and their odd correspondences]
Let $E$ be any elliptic curve over an algebraically closed field.
Every irreducible surface $V\subset E^4$ of product-polarization degree
at most $6$ is a translate of an abelian surface isomorphic to $E^2$,
and is a graph over some pair of coordinate elliptic curves.

Relative to the fixed splitting $E^4=B\times B$, such a surface either
has both odd K\"unneth correspondence components zero, or has exactly
one nonzero odd component. In the latter case it is described by a
rank-one matrix $U\in M_2(\End(E))$ with
\[
 T_U=\sum_{i,j}\deg U_{ij}\leq2,
 \qquad \deg_\theta V=2(1+T_U).
\]
Here degree zero is assigned to a zero endomorphism.
\end{lemma}

\begin{proof}
For each coordinate pair $i<j$, let $a_{ij}$ be the projection degree if
the projection $V\to E_i\times E_j$ is dominant, and zero otherwise.
Since the square of a point-divisor pulled back from one elliptic curve
is zero,
\[
 \frac{\deg_\theta V}{2}=\sum_{i<j}a_{ij}\leq3.
\]
At least one pair projection is dominant: the four coordinate maps
together have two-dimensional image, so two of the corresponding
function fields are algebraically independent. If no positive $a_{ij}$
were one, only one pair could be dominant. Every other coordinate map
would then be constant. Indeed a nonconstant third coordinate, dependent
with each coordinate of an independent pair, would make both original
coordinates algebraic over its one-dimensional function field, a
contradiction. With the other coordinates constant, $V$ is precisely the
product of that pair and has projection degree one. This contradiction
proves that a degree-one pair exists.

The other coordinate maps are rational maps from this $E^2$ to $E$.
They extend to morphisms and are homomorphisms followed by translations,
by Weil extension and rigidity. Their closed graph equals $V$. After
translation, write the two remaining maps as a matrix $h$. Its graph
degree is
\[
 2(1+\deg h+T_h),\qquad T_h=\sum\deg h_{ij}.
\]
If $h$ has rank two, then $\deg h\geq1$ and
$T_h\geq2\sqrt{\deg h}\geq2$ by the surface Hodge index inequality.
The degree budget therefore forces $\rank h\leq1$ and $T_h\leq2$.

If the input coordinate pair is one of the two fixed $B$ factors,
this is a rank-one graph or its transpose. It has one odd coefficient
$U=h$ and the other zero, unless $h=0$.

For a pair crossing the two $B$ factors, relabel coordinates within
those factors and write the two projections from the parameter $E^2$ as
\[
 f=\begin{pmatrix}1&0\\ a&b\end{pmatrix},\qquad
 g=\begin{pmatrix}0&1\\ c&d\end{pmatrix}.
\]
The remaining-coordinate matrix is
$h=\left(\begin{smallmatrix}a&b\\c&d\end{smallmatrix}\right)$, of rank
at most one and norm sum at most two. The entries $b,c$ cannot both be
nonzero: the budget would make them units and $a=d=0$, giving rank two.
If $b=c=0$, both projections have at most curve image, and $V$ is a
product of curves, with no odd component.

Suppose $b\ne0,c=0$. Then $f$ is an isogeny, $g$ has rank one, and the
unique odd coefficient, under the usual correspondence convention, is
\[
 U=(\deg b)g f^{-1}
  =\begin{pmatrix}-\bar b a&\bar b\\
                  -d\bar b a&d\bar b\end{pmatrix}.
\]
The bar denotes the dual elliptic endomorphism. If $\deg b=2$, the
budget forces $a=d=0$, so $T_U=2$. If $\deg b=1$, at most one of $a,d$
is nonzero, and it is a unit; then $T_U=1$ or $2$, equal to $T_h$.
Thus the asserted degree formula holds. The case $c\ne0,b=0$ is its
transpose. These formulas follow from $f_*g^*$ on $H^1$ and duality on
$H^3$; the other odd component vanishes because a rank-one homomorphism
has zero pullback on $H^3$. All degrees include inseparable degree.
\end{proof}

\begin{theorem}[V03: exact complexity of the primitive class]
For the product principal polarization,
\[
 c_\C(t_0)=16.
\]
The same exact value holds at every geometric good reduction of residue
characteristic $p>3$, for one prime-to-$p$ realization and for the
simultaneous-realization convention. In particular, it holds at every
ordinary good prime. The argument below makes no assertion about the
primitive minimum in characteristics two and three.
\end{theorem}

\begin{proof}
The known graph expression for $t_*$ represents $2t_0$ with degrees
$16$ on each side, so the upper bound is $16$. In a hypothetical
representation of complexity below $16$, both effective degrees are
at most $14$, since all surface degrees are even. The previously proved
degree-restricted Weil-coordinate estimates exclude every multiplier
$N\geq2$, because $|z(Nt_0)|=N\sqrt3>2\sqrt2$.

Consider $N=1$. At $p>3$, or in characteristic zero, every
nonzero-Weil irreducible component of degree at most $14$ dominates
both $B$ factors and is a translated graph of an isogeny of degree one
or two, possibly after transposing the factors. Each costs at least
$8$, so there is at most one in each numerator. Both numerators must
have such a component: after the argument below forces the main
components to be automorphism graphs, their coordinates have modulus
one, whereas $|z(t_0)|=\sqrt3$. Alternatively, before that step,
the one-component bound $|z|\leq\sqrt2<\sqrt3$ already excludes a
single nonzero-Weil numerator.

Call these the two main surfaces. Each remaining component has degree
at most $6$ and is classified by V02. A graph of a degree-$d$ isogeny
has odd correspondence coefficients $h$ and $d h^{-1}$, whose isogeny
degrees are $d$ and $d^3$; transpose interchanges the pair. For $d=1$
the degree pair is $(1,1)$, and for $d=2$ it is $(2,8)$ or $(8,2)$.
Equality of either odd coefficient for two main surfaces therefore
forces the same orientation and isogeny, or the identical graph viewed
in the opposite orientation when $d=1$. In particular it forces the
same Weil coordinate. Since their Weil-coordinate difference is
$z(t_0)\ne0$, \emph{both} odd coefficient differences are nonzero.

The target $t_0$ has no odd K\"unneth component. At least one correction
of each of the two orientations is therefore necessary. Every such
correction costs at least $4$, and each numerator has residual degree
at most $6$, so it contains at most one correction. Hence there are
exactly two corrections, one in each numerator and one of each
orientation. Both main surfaces consequently have degree at most $10$.
A degree-two isogeny graph has degree at least
$2(1+2+\lceil2\sqrt2\rceil)=12$, so both main surfaces are
automorphism graphs, denoted $\Gamma_M,\Gamma_N$.

Their entry norm sums $T_M,T_N$ are at most $3$ and at least $2$.
In characteristic zero or characteristic $p>3$, every degree-one
endomorphism of the $j=0$ elliptic curve lies in $\mathcal O_K^\times$.
A non-$K$ entry has norm at least two. Under the budget at most three,
such an entry in an invertible matrix would leave exactly one other
unit entry, producing degree two, not an automorphism. Thus $M,N$ are
matrices over $\mathcal O_K$, with either two or three unit entries and
all remaining entries zero.

Their determinants are units. Since
$z(\Gamma_M)-z(\Gamma_N)=1-\omega$, their determinant ratio is
$\omega$ or $\omega^2$: two complex units at distance $\sqrt3$ have
that ratio. The correction for the $H^1$ correspondence gives
\[
 U=\pm(M-N),\qquad \rank U=1,\qquad T_U\leq2.
\]
If either main matrix, say $M$, has norm sum two, it is a monomial
unitary matrix. Put $V=M^{-1}(M-N)$, also rank one, with the same
Frobenius norm. The rank-one determinant identity gives
\[
 |\operatorname{tr}V|^2
  =|1-\det N/\det M|^2=3.
\]
For a rank-one complex matrix $V=ab^*$,
$|\operatorname{tr}V|\leq\|a\|\|b\|=\|V\|_F$.
This contradicts $\|V\|_F^2=T_U\leq2$. The same argument applies
with $M,N$ interchanged.

It remains that both main norm sums are three, so both graph degrees
are ten and each residual degree is at most four. V02 now gives
$T_U\leq1$. The difference $M-N$ therefore changes exactly one matrix
entry by a unit. Both matrices have exactly three nonzero unit entries,
so their zero position is unchanged. Their determinant difference is
that unit times the corresponding cofactor, which is either zero or
a unit. Its squared modulus is therefore zero or one. But the
determinant ratio $\omega$ or $\omega^2$ forces squared modulus three.
This final contradiction excludes $N=1$ and proves the lower bound.

The upper bound uses actual global graph cycles and therefore holds in
all compatible realizations. The lower bound in one fixed realization
already proves the simultaneous statement.
\end{proof}

\begin{proposition}[V04: the self-intersection norm does not determine complexity]
On this one fixed complex polarized product there are rational algebraic
Weil classes of self-intersection $6$ whose least complexities are
unbounded.
\end{proposition}

\begin{proof}
Put $u=t_0$ and $R=[\omega]^*$. Their Gram matrix is
$\left(\begin{smallmatrix}6&-3\\-3&6\end{smallmatrix}\right)$.
For positive integers $n$ divisible by three, set
\[
 q_n=n^2-n+1,\qquad a_n=n^2-1,\qquad b_n=2n-1,\qquad
 v_n=\frac{a_nu+b_nRu}{q_n}.
\]
The identity $a_n^2-a_nb_n+b_n^2=q_n^2$ gives $v_n^2=6$.
Also $\gcd(q_n,a_n,b_n)=1$: a common divisor divides both $n-2$
and $2n-1$, hence divides three, while $q_n\equiv1\pmod3$.

Let $\Lambda=W\cap H^4(A,\Z(2))$ and
$\Lambda_0=\Z u+\Z Ru$. If $h=[\Lambda:\Lambda_0]$, the Gram
determinant formula gives $h^2\mid27$, so $h=1$ or $3$.
If $Nv_n$ is integral, then $hNv_n\in\Lambda_0$, forcing
$q_n\mid hN$. Since $q_n$ is prime to three, it follows that
$q_n\mid N$. Conversely $q_nv_n$ is integral, so its exact integral
cohomology denominator is $q_n$.

Every integral cycle representative of $Nv_n$ must therefore have
$N\geq q_n$, and $c_\theta(v_n)\geq q_n\to\infty$.
All the classes are rationally algebraic, since $u,Ru$ are. Explicitly,
the representatives of $2u$ and $2Ru$ give
\[
 c_\theta(v_n)\leq
 \max\bigl(2q_n,16(|a_n|+|b_n|)\bigr).
\]
Thus this unboundedness keeps the variety, polarization, and norm fixed;
the rational lattice denominators account for it.
\end{proof}

\begin{lemma}[V05: all small nonzero-Weil surfaces are abelian]
In every characteristic, an irreducible surface $V\subset E^4$ of
product-polarization degree at most $14$ with nonzero Weil projection
is a translate of an abelian surface isomorphic to $E^2$. More precisely,
one of its six coordinate-pair projections has degree one, so it is a
graph over that pair.
\end{lemma}

\begin{proof}
Let $a_{ij}$ be the six pair-projection degrees as in V02, so
$\sum_{i<j}a_{ij}\leq7$. A constant coordinate forces zero Weil
projection: each of the two Weil exterior lines uses one $H^1$ factor
from each elliptic coordinate. Thus all four coordinates are nonconstant.

For their function subfields inside $k(V)$, define $i\sim j$ if the
pair projection has curve image, or equivalently the two one-dimensional
function fields are algebraically dependent. This is an equivalence
relation. For transitivity, the composita of the fields for $i,j$ and
for $k,j$ are algebraic over the field for $j$; their combined
compositum still has transcendence degree one. Pairs in different
equivalence classes are precisely the
dominant pair projections.

If no $a_{ij}$ equals one, each positive one is at least two, so there
are at most three pairs in different classes. For four nonconstant
coordinates spanning a two-dimensional function field, the only
possible partition is $(1,3)$: a partition $(2,2)$ gives four crossing
pairs, $(2,1,1)$ gives five, and $(1,1,1,1)$ gives six, whereas a single
class would have total transcendence degree one. Therefore $V$ is
$E_i\times C$, where $C$ is the curve image in the other three factors.
Indeed it is a closed integral surface inside this integral product of
dimension two, hence equals it.

Pull the Weil forms back to $E_i\times\widetilde C$, with
$\widetilde C$ the normalization. Three of their four degree-one
factors come from the curve, so their cup product is zero since
$H^3(\widetilde C)=0$. Thus this product has zero Weil projection, a
contradiction. A degree-one pair must exist. Extension of the remaining
rational coordinate maps and rigidity, as in V02, give the claimed
abelian graph.
\end{proof}

\begin{theorem}[V06: the primitive minimum at every good reduction]
The exact equality
\[
                    c_\theta(t_*/2)=16
\]
holds at every geometric good reduction, including characteristics two
and three, for one fixed prime-to-characteristic realization and for
the simultaneous-realization convention.
\end{theorem}

\begin{proof}
Only characteristics two and three remain after V03. The degree-$16$
upper bound is already supplied by the global representative of
$2t_0=t_*$. The previously proved bounds for an effective cycle of degree
at most $14$ are $|z|\leq5/3$ in characteristic two and
$|z|\leq\sqrt2$ in characteristic three. Therefore
$|z(Nt_0)|=N\sqrt3$ excludes every $N\geq2$.

Suppose that $t_0=[Z^+]-[Z^-]$ with both effective degrees at most $14$.
Each numerator has exactly one irreducible component with nonzero Weil
projection, of multiplicity one. Existence on both sides follows from
the same one-cycle bounds, which are strictly less than $\sqrt3$;
uniqueness follows because each nonzero-Weil component costs at least
$8$. Call these two components the main surfaces. All remaining
components have degree at most $6$ and are classified by V02. In
particular each numerator contains at most one component with nonzero
odd correspondence, and its odd matrix has rank one and norm sum at
most two. If the main degree is at least ten, the residual budget is
at most four, so that odd matrix has at most one nonzero entry, a unit.

By V05 each main surface is a translated abelian surface. According to
its two projections onto $B$, it has one of the following types:
\begin{itemize}
\item $\mathsf D$: both projections are dominant;
\item $\mathsf L$ or $\mathsf R$: exactly one projection is dominant,
and the other has curve image;
\item $\mathsf P$: both projections have curve image, so the surface
is a product of two curves.
\end{itemize}
A point image would give zero Weil projection and does not occur.
We exclude every pair of types.

\paragraph{Two dominant main surfaces.}
The odd-degree-pair argument of V03 forces one correction of each odd
orientation, and hence both main degrees at most ten. Thus they are
automorphism graphs $\Gamma_M,\Gamma_N$ with entry norm sums two or
three. Every nonzero entry is a unit: a norm-two entry would leave only
one other unit and produce a degree-two isogeny. This assertion does
not require the units to lie in $K$. Let $\pi_K$ be the positive-norm
orthogonal projection of the quaternion algebra onto $K$, and recall
\[
 |z(\Gamma_M)-z(\Gamma_N)|
   =|\det(\pi_K M)-\det(\pi_K N)|=\sqrt3.
\]
The $H^1$ correction gives $R=M-N$ of rank one and norm sum at most
two. If both main norm sums are three, that bound improves to one.

If both have norm sum two, they are monomial matrices. Different
monomial supports would give four unit entries in $R$, impossible.
With the same support, rank one forces one common entry. The determinant
difference of the $K$ projections then has modulus at most
$\sqrt{T_R}\leq\sqrt2$, since projection decreases norm.

If one main norm sum is two and the other three, permute rows and
columns to make the first matrix diagonal. If the second matrix lacks
a diagonal entry, their difference has at least three unit entries,
impossible. Otherwise the second is triangular. Rank one of the
difference forces a common diagonal entry. Its remaining diagonal
difference has norm at most one, because the off-diagonal unit already
costs one. The $K$-determinant difference has modulus at most one.

Finally, if both norm sums are three, the norm-one difference changes
one entry. Both matrices have three nonzero entries, so the zero
position is unchanged. The $K$-determinant difference is the $K$
projection of that entry difference times the $K$ projection of a
common unit cofactor, and again has modulus at most one. Every case
contradicts the required modulus $\sqrt3$.

\paragraph{One dominant and one nondominant main surface.}
A dominant main has two full-rank odd matrices. A nondominant main
has at least one zero odd matrix, and its other odd matrix has rank
at most one. In the block where it is zero, at least two rank-one
corrections are needed to cancel the dominant main's full-rank matrix.
The other odd difference is also nonzero and needs a third correction.
Only two corrections fit in total. Thus a pair involving exactly one
$\mathsf D$ surface is impossible.

\paragraph{Numerical and support facts for $\mathsf L$ and $\mathsf R$.}
Treat $\mathsf L$; the other type is its transpose. Translate the
surface through the origin, and let its projections be an isogeny
$f:V\to B$ of degree $a$ and a rank-one map $g:V\to B$. Put
$q=g f^{-1}\in\Hom^0(B,B)$. Its single nonzero odd coefficient is
$R=a q$. If $c=\lambda_1\lambda_2[V]$, projection formula gives
\[
 c=a\sum_{i,j}\operatorname{Nm}(q_{ij}),\qquad
 \sum_{i,j}\operatorname{Nm}(R_{ij})=ac.
\]
The middle K\"unneth component factors as the image curve class
times the effective generic-fiber divisor class, as in C10/F03.
Nonzero Weil projection forces both curve degrees at least two.
Consequently $c$ is four or six, and the only possibilities are
\[
\begin{array}{c|c|c|c}
 \deg_\theta V&a&c&T_R\\ \hline
 10&1&4&4\\
 12&2&4&8\\
 14&3&4&12\\
 14&1&6&6
\end{array}
\]
Both the image curve and the kernel curve have nonzero projections to
both elliptic coordinates: a coordinate curve has zero nontrivial
scalar-$\omega$ part. Writing the rank-one matrix $R$ as a column
times a row therefore shows that all four entries of $R$ are nonzero.

\paragraph{An $\mathsf L$ main paired with $\mathsf P$ or $\mathsf R$.}
Against a product main, its four-entry odd matrix must be canceled
by the remaining components. The $\mathsf L$ side has residual degree
at most four and supplies at most one matrix entry; the product side
has residual degree at most six and supplies at most two. Three
positions cannot cancel four nonzero entries. Against the opposite
orientation $\mathsf R$, each main has its four-entry matrix in a
different odd block, and both residual budgets are at most four;
the same support obstruction is stronger. The transposed cases are
identical.

\paragraph{Two main surfaces of the same one-dominant type.}
Let their rank-one four-entry matrices be $R,R'$. If either main degree
is at least twelve, that numerator has no odd correction. The other
can change at most one entry. Two rank-one matrices over a division
algebra, all of whose entries are nonzero, cannot differ in exactly
one entry: any three entries determine the fourth. Hence $R=R'$.
The four distinct norm sums in the table determine the same pair
$(a,c)$, so the rational graph maps $q=R/a$ agree. A rational graph
map determines a unique connected abelian subgroup of $B\times B$:
for an integer clearing its denominator, take the image of
$x\mapsto(nx,nqx)$. Thus the main surfaces have the same cohomology
class, contradicting their nonzero Weil difference.

The remaining case has both main degrees ten. Then $a=1$ and the
surfaces are literal rank-one graphs $h,h'$ with four unit entries.
Their six coordinate-pair projection degrees are identical: one for
the first $B$, zero for the second, and one for each cross pair.
If the main odd difference is nonzero, both corrections must lie in
that same odd orientation: a lone correction in the other orientation
could not cancel, while one correction alone cannot change a four-entry
rank-one matrix in only one position.
Each correction, if present, has degree four and consists of a
rank-one graph with exactly one nonzero unit entry. One correction
alone, or two in the same matrix position, would make $h-h'$ supported
in a single entry and force equality as before. If the two corrections
occupy distinct positions, their coordinate-pair projection-degree
vectors differ. For a single-entry graph, that vector has a one for
the first $B$ pair and a one for the unique cross pair complementary
to the chosen input coordinate; it identifies the matrix position.
There is no residual degree left for other components. Equality to
the pure Weil class requires equality of all six coordinate-pair
degrees, so distinct correction positions are also impossible.

\paragraph{Two product main surfaces.}
Write them as $C^+\times D^+$ and $C^-\times D^-$. Their two curve
degrees are $(2,2)$, $(2,3)$, or $(3,2)$, so their surface degrees are
eight or twelve. Decompose
\[
 H^2(B)=I\oplus J\oplus T,
\]
where $I$ is spanned by the two coordinate-fiber classes, $J$ is the
scalar-$\omega$-invariant cross part, and $T$ is the sum of the two
nontrivial scalar-$\omega$ eigenlines. The Weil plane is contained in
$T\otimes T$. Its nonzero rational vectors, in particular $t_0$,
have tensor rank two: over a splitting field their two conjugate
opposite-character coefficients are both nonzero.

Every remaining component of degree at most six has middle class in
$I\otimes H^2(B)+H^2(B)\otimes I$. Indeed the curve-image
factorization has mixed degree at most three, so at least one of its
two effective curve factors has degree one and is a coordinate fiber.
This applies also to components having an odd correspondence.

The rank-two target in $T\otimes T$ implies that the two $T$ parts of
$C^+,C^-$ are linearly independent, and likewise those of $D^+,D^-$.
Projection of the full equality to $T\otimes J$ therefore forces
both $J$ parts of $D^+,D^-$ to vanish. Projection to $J\otimes T$
does the same for $C^+,C^-$. Each of these four curves is a graph or
transposed graph of a norm-one or norm-two elliptic endomorphism,
and its $J$ part vanishes exactly when the endomorphism's $K$ part
vanishes.

In characteristic two, no Hurwitz unit has zero $K$ projection:
its squared projected norm is one or $1/3$. Each product main has at
least one degree-two curve, hence a unit graph. This is a contradiction.

In characteristic three, a norm-two endomorphism has projected norm
one, whereas a non-$K$ unit has projected norm zero. Thus all four
curves must be non-$K$ unit graphs, both main surfaces have degree
eight, and the $I$ part of every curve is
$\lambda=\theta_1+\theta_2$. The unwanted mixed components are
\[
 (T(C^+)-T(C^-))\otimes\lambda,
 \qquad \lambda\otimes(T(D^+)-T(D^-)),
\]
and both are nonzero by the independence just proved. A degree-at-most-six
component with a nonzero $T$ contribution costs at least four and
has a single coordinate-fiber factor on the other side; it contributes
to at most one of $T\otimes I$ and $I\otimes T$. At least two
such components are needed for each displayed direction, one for
each of the two independent coordinate-fiber classes. At least four
are therefore needed in total, costing at least sixteen. The two
residual budgets together are at most twelve. This final contradiction
excludes the product pair.

All pairs of main types have been excluded, so no multiplier-one
representation with degrees at most fourteen exists. Together with
the earlier exclusion of larger multipliers, parity of degrees, and
the global upper bound, this proves the exact value sixteen.
\end{proof}
